\capitulo{1}{Introducción}

El modelado tridimensional ha adquirido un papel esencial en diversos campos como videojuegos, animación, diseño industrial, arquitectura o visualización científica. En este ámbito, las herramientas de modelado 3D se encargan de crear, modificar y examinar modelos con una amplia selección de métodos y flujo de trabajo~\cite{lavioila2017}.

A pesar de esta creciente popularidad, el proceso de modelado tridimensional sigue siendo complejo e intrínsecamente ligado a la experiencia de quien lo realiza. Varias investigaciones han revelado que la interacción en espacios tridimensionales tiene ciertas desventajas relacionadas con manipular el espacio, comprender la geometría y realizar varias tareas al mismo tiempo~\cite{jankowski2014}.

Otras plataformas extendidas en la industria, como Blender, no tienen mecanismos nativos de registros estructurados de la interacción del usuario ni métodos que permitan observar el estado interno de la escena durante el diseño de la misma. Esto dificulta obtener métricas objetivas, hacer un análisis sistemático del flujo de trabajo, optimizar procesos, automatizar tareas y evaluar el rendimiento.

La comprensión de cómo interactúan los usuarios con entornos digitales se considera un tema fundamental en el campo de la persona-computadora-persona, especialmente con respecto a los sistemas tridimensionales~\cite{lavioila2017}. Recopilar información sobre lo que los usuarios hacen ayuda a determinar los patrones de usabilidad, las operaciones habituales y las técnicas de modelado. Los últimos años han visto propuestas variadas para capturar y analizar los datos en los entornos 3D, incluida la realidad mixta e herramientas inteligentes~\cite{3description2024,malik2023}. Sin embargo, todavía existe una brecha metodológica en la vigilancia analítica no invasiva dentro de los entornos utilizados habitualmente.

La motivación principal de la presente investigación está en resolver dicha deficiencia, proporcionando un marco analítico para que tanto instituciones educativas como profesionales evalúen la capacidad y la habilidad técnica en tareas de modelado. De ahí que la recolección de datos resulta valiosa como fundamento de métodos de aprendizaje supervisado, mejoramiento de interfaces de uso y auditoría de flujos. 

Dicho contexto hace surgir la necesidad de crear soluciones de software capaces de registrar y analizar los eventos en tiempo real. De ahí que el objetivo principal del trabajo presente consista en elaborar y implantar un sistema automatizado que reciba, estruture y analice los datos durante las sesiones de modelado en Blender. Para esto se propone el desarrollo de un sistema modular, basado en \textit{add-ons}, cuya tarea consiste en registrar las acciones de modelado, junto con variables structurales del entorno, es decir, el estado de la escena, la topología de los objetos, y la configuración de los modificadores.

La solución desarrollada se estructura en dos complementos diferenciados. El primero actúa como sistema de captura y registra de forma no intrusiva la actividad del usuario en ficheros CSV. El segundo procesa estos registros, calcula métricas estadísticas y geométricas, y presenta los resultados mediante una interfaz integrada en Blender. De este modo, se proporciona una plataforma funcional que facilita el diagnóstico del proceso de modelado y abre la posibilidad de establecer futuras líneas de optimización o automatización en entornos de diseño tridimensional.

\section{Estructura de este trabajo}

Este documento se organiza en siete capítulos principales y una colección de anexos técnicos. La memoria mantiene una descripción sintética del proyecto con el fin de no sobrecargar el cuerpo principal del documento, mientras que los anexos recogen el detalle formal de requisitos, diseño, programación, instalación y uso.

\begin{itemize}

\item \textbf{Capítulo 1. Introducción}\\
Presenta el contexto general del modelado tridimensional, la problemática identificada y la solución propuesta.

\item \textbf{Capítulo 2. Objetivos del proyecto}\\
Define el objetivo general y los objetivos específicos relacionados con la captura, el análisis, la visualización y la validación del sistema.

\item \textbf{Capítulo 3. Conceptos teóricos}\\
Resume los fundamentos necesarios para comprender el trabajo: Blender, \textit{add-ons}, monitorización de eventos, estructuras \texttt{Mesh} y \texttt{BMesh}, métricas geométricas y visualización de datos.

\item \textbf{Capítulo 4. Técnicas y herramientas}\\
Describe la metodología de desarrollo y las tecnologías utilizadas, incluyendo Python, \texttt{bpy}, \texttt{bmesh}, NumPy, Matplotlib, CSV y herramientas de control de versiones.

\item \textbf{Capítulo 5. Aspectos importantes del desarrollo del proyecto}\\
Explica la arquitectura de la solución, el funcionamiento de los dos \textit{add-ons}, el flujo de datos, las métricas implementadas, las decisiones técnicas, los problemas encontrados y la validación realizada.

\item \textbf{Capítulo 6. Trabajos relacionados}\\
Sitúa la solución desarrollada en relación con investigaciones previas sobre análisis de interacción, registros de eventos, modelado 3D y aplicaciones educativas.

\item \textbf{Capítulo 7. Conclusiones y trabajo futuro}\\
Recoge las conclusiones generales y técnicas del proyecto, las limitaciones detectadas y las líneas de ampliación propuestas.

\item \textbf{Anexos}\\
Incluyen la documentación complementaria: plan de proyecto, especificación de requisitos, diseño, documentación técnica de programación, manual de usuario y sostenibilización curricular.

\end{itemize}
